\section{Related Work}
\label{sec:related}

\paragraph{Action recognition and the acrobatic gap.}
Large-scale action recognition benchmarks such as Kinetics-700~\cite{kay2017kinetics,carreira2019shortkinetics}
and NTU-RGBD~\cite{shahroudy2016ntu} have driven rapid progress in generic
video classification, but their label sets are dominated by daily activities
and contain essentially no acrobatic motion. AMASS~\cite{mahmood2019amass}, the
de-facto training corpus for human mesh recovery, has the same blind spot:
motion-capture sequences of walking, dancing, and sitting, with no flips,
inversions, or multi-twist rotations. Models trained on these distributions
transfer poorly to parkour and gymnastics footage.

\paragraph{Human mesh recovery for sports.}
GVHMR~\cite{shen2024gvhmr}, WHAM~\cite{shin2024wham}, TRAM~\cite{wang2024tram},
and TokenHMR~\cite{dwivedi2024tokenhmr} estimate SMPL~\cite{loper2015smpl}
bodies in world coordinates from monocular video. GVHMR reports
gravity-aligned trajectories, which in principle should allow clean
flip-and-twist decomposition from the global orientation. In practice, all of
these models inherit AMASS's distribution gap and produce unreliable
rotations on acrobatic motion~--- a finding we confirm empirically in
Sec.~\ref{sec:journey}. AthletePose3D~\cite{yeung2025athletepose3d} attacks
the problem directly by fine-tuning on extreme poses, but its target
disciplines are diving and gymnastics, not parkour, and public checkpoints
are limited.

\paragraph{Vision--language models as zero-shot recognizers.}
A recent line of work uses general-purpose vision--language models
(VLMs)~--- Flamingo~\cite{alayrac2022flamingo}, GPT-4V~\cite{openai2023gpt4v},
Claude~\cite{anthropic2024claude3}, Gemini~\cite{team2024gemini},
Qwen-VL~\cite{bai2023qwen}~--- as zero-shot recognizers for tasks with no
labeled training set. These models inherit broad world knowledge from
internet-scale pretraining; several of them demonstrably recognize FIG parkour
trick names from a handful of frames. Our contribution is not a new VLM, but
a pragmatic system that turns this raw naming capability into a
FIG-compliant judge through segmentation, a fuzzy trick matcher, and a
D-score assembly pipeline.

\paragraph{Automated judging in gymnastics and related sports.}
Fujitsu's Judging Support System~\cite{fujitsu2023jss} is the most visible
deployed example: a multi-camera 3D setup that tracks joints with custom
hardware and reports element-level scores to FIG judges for artistic
gymnastics. The ISU has published similar pilots for figure
skating~\cite{isu2023automation}. These systems achieve high precision but
require bespoke capture rigs, closed datasets, and per-element engineering.
No open-source equivalent exists for parkour, and the FIG Parkour Code of
Points~\cite{fig2025codepoints} itself is only three years old. \sys{} targets
the monocular smartphone-footage regime that parkour competitions actually
use, trading precision for deployability.
